\begin{table}[t]
\centering
\small
\setlength{\tabcolsep}{4pt}
\caption{Lightweight evaluation on the public PSS/HyperPlonk artifact. The 8-party column exercises representative collaborative release objects; the 16-party point is a commitment-only scalability sanity check.}
\label{tab:second-inst-light}
\begin{tabular}{lcc}
\toprule
Check & 8 parties & 16 parties \\
\midrule
Sumcheck released rounds vs. reference & $8/8$ & -- \\
Modified Sumcheck round & WITHHOLD & -- \\
Packed PCS commitment & MATCH & MATCH \\
PCS opening proof & MATCH & -- \\
Native PCS verification$^\dagger$ & PASS & -- \\
Context/ordinal/payload binding & PASS & -- \\
\bottomrule
\end{tabular}
\vspace{1mm}
\begin{minipage}{0.96\columnwidth}
\footnotesize
$^\dagger$The upstream benchmark keeps the final scalar returned by \texttt{c\_open} in an internal share representation. We therefore verify the collaboratively reconstructed commitment and opening proof using the corresponding upstream reference opening value; we make no correctness claim for that benchmark-only scalar reconstruction path.
\end{minipage}
\end{table}

\begin{table}[t]
\centering
\small
\caption{Per-release diagnostic costs for the 8-party second-instantiation experiment (median of five measured runs after one warm-up). These numbers are not full-system PVIA overhead.}
\label{tab:second-inst-cost}
\begin{tabular}{lrr}
\toprule
Operation & Protocol path & Release check \\
\midrule
Collaborative Sumcheck & 7.29 ms & 67.3 $\mu$s \\
Collaborative PCS commit/open & 372.52 ms & 30.76 ms \\
Context binding & -- & 0.80 $\mu$s \\
\bottomrule
\end{tabular}
\end{table}
